% GENERATED FILE. Edit canonical lessons, then run npm run generate:books.
% canonical-source-hash: fnv1a64:6d8595d9663931e4
% canonical-lessons: TE-C475-sangamam, TE-C475-anakatta, TE-C475-jalasayam, TE-C475-nitiparudala, TE-C475-cittadi

\chapter{Rivers and Dams}
\label{ch:te-rivers-and-dams}
\hlchaptermodality{475}

\begin{chapteropening}
\textbf{By the end of this chapter:} \emph{I can say a confluence, a dam, a reservoir, irrigation and a marsh.}

Complete the last lesson of chapter 475: I can say a confluence, a dam, a reservoir, irrigation and a marsh.
\end{chapteropening}

\section[\texorpdfstring{\te{సంగమం} (saṅgamaṁ)}{సంగమం (saṅgamaṁ)}]{\te{సంగమం} (saṅgamaṁ) — a confluence}

\label{lesson:TE-C475-sangamam}



\begin{quote}
Before the new one: say the Telugu for a volcano, then the Telugu for peaceful.
\end{quote}

\subsection*{You'll want to know: \te{సంగమం}}
\textbf{\te{సంగమం}} — \emph{saṅgamaṁ} — \textquotedblleft{}a confluence\textquotedblright{}.

\begin{quote}
\te{రెండు} \te{నదుల} \te{సంగమం} \te{దగ్గర} \te{పెద్ద} \te{జాతర} \te{జరుగుతుంది}.
\end{quote}

\begin{quote}
\emph{reṇḍu nadula saṅgamaṁ daggara pedda jātara jarugutundi.}
\end{quote}

\begin{quote}
A big fair is held where the two rivers meet.
\end{quote}

The place where two rivers meet; also any meeting of two things.

\begin{cousinweb}
\textbf{\te{సంగమం}} (\emph{saṅgamaṁ}) is borrowed from Sanskrit \emph{saṅgama}, \textquotedblleft{}a coming together, a confluence of rivers\textquotedblright{} (Monier-Williams), from \emph{sam-gam}, \textquotedblleft{}to go together\textquotedblright{}.
\end{cousinweb}

\begin{grammarlens}[title={What you've built}]
One more word for getting around.
\end{grammarlens}

\subsection*{Your turn}
\begin{itemize}
\raggedright
  \item \emph{Say it:} \emph{saṅgamaṁ}
  \item \emph{Say it:} \emph{saṅgamaṁ}, once more
  \item \emph{Say it:} \emph{reṇḍu nadula saṅgamaṁ daggara pedda jātara jarugutundi.} — A big fair is held where the two rivers meet.
  \item \emph{Recall:} say the Telugu for a volcano, then the Telugu for peaceful, then say \emph{saṅgamaṁ} again
\end{itemize}

\begin{tcolorbox}[breakable,colback=teal!4,colframe=teal!35!black,title={Before you move on}]
What does \te{సంగమం} mean? (\textquotedblleft{}A confluence\textquotedblright{}.) Say it once more.
\end{tcolorbox}
\section[\texorpdfstring{\te{ఆనకట్ట} (ānakaṭṭa)}{ఆనకట్ట (ānakaṭṭa)}]{\te{ఆనకట్ట} (ānakaṭṭa) — a dam; a weir}

\label{lesson:TE-C475-anakatta}



\begin{quote}
Before the new one: say the Telugu for peaceful, then the Telugu for a confluence.
\end{quote}

\subsection*{You'll want to know: \te{ఆనకట్ట}}
\textbf{\te{ఆనకట్ట}} — \emph{ānakaṭṭa} — \textquotedblleft{}a dam; a weir\textquotedblright{}.

\begin{quote}
\te{నది} \te{మీద} \te{కొత్త} \te{ఆనకట్ట} \te{కట్టారు}.
\end{quote}

\begin{quote}
\emph{nadi mīda kotta ānakaṭṭa kaṭṭāru.}
\end{quote}

\begin{quote}
They built a new dam on the river.
\end{quote}

\te{ఆన}, a bank, + \te{కట్ట}, a bund. \te{ఆనకట్ట} \te{కట్టు}, to build a dam.

\begin{cousinweb}
\textbf{\te{ఆనకట్ట}} (\emph{ānakaṭṭa}) is native Dravidian: \te{ఆన}, \textquotedblleft{}a bank, a dam\textquotedblright{}, which DEDR 122 lists as Telugu \emph{ana} with a query, and \te{కట్ట}, \textquotedblleft{}a bund, an embankment\textquotedblright{}, from \te{కట్టు}, \textquotedblleft{}to tie, to build\textquotedblright{} (DEDR 1147).
\end{cousinweb}

\begin{grammarlens}[title={What you've built}]
One more word for getting around.
\end{grammarlens}

\subsection*{Your turn}
\begin{itemize}
\raggedright
  \item \emph{Say it:} \emph{ānakaṭṭa}
  \item \emph{Say it:} \emph{ānakaṭṭa}, once more
  \item \emph{Say it:} \emph{nadi mīda kotta ānakaṭṭa kaṭṭāru.} — They built a new dam on the river.
  \item \emph{Recall:} say the Telugu for peaceful, then the Telugu for a confluence, then say \emph{ānakaṭṭa} again
\end{itemize}

\begin{tcolorbox}[breakable,colback=teal!4,colframe=teal!35!black,title={Before you move on}]
What does \te{ఆనకట్ట} mean? (\textquotedblleft{}A dam; a weir\textquotedblright{}.) Say it once more.
\end{tcolorbox}
\section[\texorpdfstring{\te{జలాశయం} (jalāśayaṁ)}{జలాశయం (jalāśayaṁ)}]{\te{జలాశయం} (jalāśayaṁ) — a reservoir}

\label{lesson:TE-C475-jalasayam}



\begin{quote}
Before the new one: say the Telugu for a confluence, then the Telugu for a dam.
\end{quote}

\subsection*{You'll want to know: \te{జలాశయం}}
\textbf{\te{జలాశయం}} — \emph{jalāśayaṁ} — \textquotedblleft{}a reservoir\textquotedblright{}.

\begin{quote}
\te{వర్షాలకు} \te{జలాశయం} \te{పూర్తిగా} \te{నిండింది}.
\end{quote}

\begin{quote}
\emph{varṣālaku jalāśayaṁ pūrtigā niṇḍindi.}
\end{quote}

\begin{quote}
With the rains the reservoir filled up completely.
\end{quote}

\te{జల}-, water, + -\te{ఆశయం}, a resting place: a place where water rests.

\begin{cousinweb}
\textbf{\te{జలాశయం}} (\emph{jalāśayaṁ}) is a compound of two words borrowed from Sanskrit: \emph{jala}, \textquotedblleft{}water\textquotedblright{}, and \emph{āśaya}, \textquotedblleft{}a resting place, a receptacle\textquotedblright{} (Monier-Williams).
\end{cousinweb}

\begin{grammarlens}[title={What you've built}]
One more word for getting around.
\end{grammarlens}

\subsection*{Your turn}
\begin{itemize}
\raggedright
  \item \emph{Say it:} \emph{jalāśayaṁ}
  \item \emph{Say it:} \emph{jalāśayaṁ}, once more
  \item \emph{Say it:} \emph{varṣālaku jalāśayaṁ pūrtigā niṇḍindi.} — With the rains the reservoir filled up completely.
  \item \emph{Recall:} say the Telugu for a confluence, then the Telugu for a dam, then say \emph{jalāśayaṁ} again
\end{itemize}

\begin{tcolorbox}[breakable,colback=teal!4,colframe=teal!35!black,title={Before you move on}]
What does \te{జలాశయం} mean? (\textquotedblleft{}A reservoir\textquotedblright{}.) Say it once more.
\end{tcolorbox}
\section[\texorpdfstring{\te{నీటిపారుదల} (nīṭipārudala)}{నీటిపారుదల (nīṭipārudala)}]{\te{నీటిపారుదల} (nīṭipārudala) — irrigation}

\label{lesson:TE-C475-nitiparudala}



\begin{quote}
Before the new one: say the Telugu for a dam, then the Telugu for a reservoir.
\end{quote}

\subsection*{You'll want to know: \te{నీటిపారుదల}}
\textbf{\te{నీటిపారుదల}} — \emph{nīṭipārudala} — \textquotedblleft{}irrigation\textquotedblright{}.

\begin{quote}
\te{ఆ} \te{ఆనకట్ట} \te{వల్ల} \te{నీటిపారుదల} \te{సులభమైంది}.
\end{quote}

\begin{quote}
\emph{ā ānakaṭṭa valla nīṭipārudala sulabhamaindi.}
\end{quote}

\begin{quote}
Thanks to that dam, irrigation became easy.
\end{quote}

\te{నీటి}, of water, + \te{పారుదల}, flowing: making water flow to the fields.

\begin{cousinweb}
\textbf{\te{నీటిపారుదల}} (\emph{nīṭipārudala}) is native Dravidian: \te{నీటి}, from \te{నీరు}, \textquotedblleft{}water\textquotedblright{} (DEDR 3690), and \te{పారుదల}, \textquotedblleft{}flowing\textquotedblright{}, from \te{పారు}, \textquotedblleft{}to run, to flow\textquotedblright{} (DEDR 3963).
\end{cousinweb}

\begin{grammarlens}[title={What you've built}]
One more word for getting around.
\end{grammarlens}

\subsection*{Your turn}
\begin{itemize}
\raggedright
  \item \emph{Say it:} \emph{nīṭipārudala}
  \item \emph{Say it:} \emph{nīṭipārudala}, once more
  \item \emph{Say it:} \emph{ā ānakaṭṭa valla nīṭipārudala sulabhamaindi.} — Thanks to that dam, irrigation became easy.
  \item \emph{Recall:} say the Telugu for a dam, then the Telugu for a reservoir, then say \emph{nīṭipārudala} again
\end{itemize}

\begin{tcolorbox}[breakable,colback=teal!4,colframe=teal!35!black,title={Before you move on}]
What does \te{నీటిపారుదల} mean? (\textquotedblleft{}Irrigation\textquotedblright{}.) Say it once more.
\end{tcolorbox}
\section[\texorpdfstring{\te{చిత్తడి} (cittaḍi)}{చిత్తడి (cittaḍi)}]{\te{చిత్తడి} (cittaḍi) — a marsh; wet ground}

\label{lesson:TE-C475-cittadi}



\begin{quote}
Before the new one: say the Telugu for a reservoir, then the Telugu for irrigation.
\end{quote}

\subsection*{You'll want to know: \te{చిత్తడి}}
\textbf{\te{చిత్తడి}} — \emph{cittaḍi} — \textquotedblleft{}a marsh; wet ground\textquotedblright{}.

\begin{quote}
\te{నది} \te{ఒడ్డున} \te{చిత్తడి} \te{నేల} \te{ఉంది}.
\end{quote}

\begin{quote}
\emph{nadi oḍḍuna cittaḍi nēla undi.}
\end{quote}

\begin{quote}
There is marshy ground on the riverbank.
\end{quote}

\te{చిత్తడి} \te{నేల}, a marsh, a wetland. As an adjective, soggy: \te{చిత్తడిగా} \te{ఉంది}.

\begin{cousinweb}
\textbf{\te{చిత్తడి}} (\emph{cittaḍi}) is native Dravidian: DEDR 2525 lists Telugu \emph{cittaḍi}, \textquotedblleft{}damp, moisture; wet\textquotedblright{}.
\end{cousinweb}

\begin{grammarlens}[title={What you've built}]
One more word for getting around.
\end{grammarlens}

\subsection*{Your turn}
\begin{itemize}
\raggedright
  \item \emph{Say it:} \emph{cittaḍi}
  \item \emph{Say it:} \emph{cittaḍi}, once more
  \item \emph{Say it:} \emph{nadi oḍḍuna cittaḍi nēla undi.} — There is marshy ground on the riverbank.
  \item \emph{Recall:} say the Telugu for a reservoir, then the Telugu for irrigation, then say \emph{cittaḍi} again
\end{itemize}

\begin{tcolorbox}[breakable,colback=teal!4,colframe=teal!35!black,title={Before you move on}]
What does \te{చిత్తడి} mean? (\textquotedblleft{}A marsh; wet ground\textquotedblright{}.) Say it once more.
\end{tcolorbox}
